\chapter{$\mathbb{R}^n$ Space}
Topology is a mathematical discipline that studies the properties of geometric objects that are preserved under continuous transformations (i.e. Studying the object without making "holes" or "destroying" the object). 

\section{$\mathbb{R}^n$ is a Normed Vector Space}
\begin{definition}[$\mathbb{R}^n$ Space]
    $\Rn$ is the set of all ordered n-tuples of real numbers. A point (or element) $\bar{x} \in \Rn$ is denoted as:
    \[
        \bar{x} = (x_1, x_2, \ldots, x_n) = \begin{bmatrix}
            x_1 \\
            x_2 \\
            \vdots \\
            x_n
        \end{bmatrix}
    \]
\end{definition}
Let $\bar{x} = (x_1, x_2, \ldots, x_n)$, $\bar{y} = (y_1, y_2, \ldots, y_n)$ and $\lambda \in \Rn$. We give $\Rn$ two operations:
\begin{citemize}
    \item Addition: $\bar{x} + \bar{y} = (x_1 + y_1, x_2 + y_2, \ldots, x_n + y_n)$
    \item Scalar multiplication: $\lambda \bar{x} = (\lambda x_1, \lambda x_2, \ldots, \lambda x_n)$
\end{citemize}
These operations satisfy the following properties for all $\bar{x}, \bar{y} \in \Rn$ and $\lambda, \lambda_1, \lambda_2 \in \R$:
\begin{citemize}
    \item $(\lambda_1 \lambda_2) \bar{x} = \lambda_1(\lambda_2 \bar{x}) = \lambda_2 (\lambda_1 \bar{x})$
    \item $1 \cdot \bar{x} = \bar{x}$
    \item $(\lambda_1 + \lambda_2) \bar{x} = \lambda_1 \bar{x} + \lambda_2 \bar{x}$
    \item $\lambda (\bar{x} + \bar{y}) = \lambda \bar{x} + \lambda \bar{y}$
\end{citemize}
There exists a basis of the vector space $\Rn$ called the standard basis, which consists of the vectors:
\[
    \bar{e}_1 = (1, 0, \ldots, 0), \quad \bar{e}_2 = (0, 1, \ldots, 0), \quad \ldots, \quad \bar{e}_n = (0, 0, \ldots, 1)
\]
Any vector $\bar{x} \in \Rn$ can be expressed as a linear combination of the standard basis vectors:
\[
    \sum_{i=1}^n x_i \bar{e}_i = (x_1, x_2, \ldots, x_n)
\]

\subsection{Scalar Product, Euclidean Norm and Cauchy-Schwarz}
\begin{definition}[Scalar Product]
    The scalar product (or dot product) of two vectors $\bar{x}, \bar{y} \in \Rn$ is defined as:
    \[
        \langle \bar{x}, \bar{y} \rangle = \sum_{i=1}^n x_i y_i = x_1 y_1 + x_2 y_2 + \ldots + x_n y_n
    \]
\end{definition}

\begin{definition}[Euclidean Norm]
    The Euclidean norm (or length) of a vector $\bar{x} \in \Rn$ is defined as:
    \[
        \|\bar{x}\| = \sqrt{\langle \bar{x}, \bar{x} \rangle} = \sqrt{x_1^2 + x_2^2 + \ldots + x_n^2}
    \]
\end{definition}
We can extract the following properties of the Euclidean norm:
\begin{citemize}
    \item $\|\bar{x} \| \geq 0$ and $\|\bar{x} \| = 0$ if and only if $\bar{x} = \bar{0} = (0, 0, \ldots, 0)$ (the zero vector)
    \item $\|\lambda \bar{x} \| = |\lambda| \|\bar{x} \|$ for all $\lambda \in \R$ and $\bar{x} \in \Rn$
\end{citemize}
Thus $\Rn$ is a normed vector space.

\begin{theorem}[Cauchy-Schwarz Inequality]
    For all $\bar{x}, \bar{y} \in \Rn$, we have:
    \[
        |\langle \bar{x}, \bar{y} \rangle| \leq \|\bar{x}\| \cdot \|\bar{y}\|
    \]
\end{theorem}
\begin{proof}
    Let $\bar{x}, \bar{y} \in \Rn$ and $\lambda \in \R$. We consider the sum:
    \[
        S = \sum_{i=1}^n (\lambda x_i + y_i)^2
    \]
    We have $S \geq 0$. Expanding $S$, we get:
    \[
        0 \leq S = \lambda^2 \sum_{i=1}^n x_i^2 + 2\lambda \sum_{i=1}^n x_i y_i + \sum_{i=1}^n y_i^2 = a\lambda^2 + b\lambda + c
    \]
    Since $S \geq 0$ for all $\lambda \in \R$, the discriminant of the quadratic polynomial $a\lambda^2 + b\lambda + c$ must be non-positive:
    \[
        b^2 - 4ac = 4 \left(\sum_{i = 1}^{n} x_i y_i\right)^2 - 4 \left(\sum_{i = 1}^{n} x_i^2\right) \left(\sum_{i = 1}^{n} y_i^2\right) \leq 0
    \]
    Which implies:
    \[
        4 \|\bar{x}\|^2 \|\bar{y}\|^2 \geq 4 \langle \bar{x}, \bar{y} \rangle^2
    \]
    Hence:
    \[
        |\langle \bar{x}, \bar{y} \rangle| \leq \|\bar{x}\| \cdot \|\bar{y}\|
    \]
\end{proof}

\begin{theorem}[Triangle Inequalities]
    For all $\bar{x}, \bar{y} \in \Rn$, we have:
    \[
        \|\bar{x} + \bar{y}\| \leq \|\bar{x}\| + \|\bar{y}\|
    \]
    and \[
        |\|\bar{x}\| - \|\bar{y}\|| \leq \|\bar{x} - \bar{y}\|
    \]
\end{theorem}
\begin{proof}
    Let $\bar{x}, \bar{y} \in \Rn$. We have:
    \[
        \|\bar{x} + \bar{y}\|^2 = \langle \bar{x} + \bar{y}, \bar{x} + \bar{y} \rangle = \|\bar{x}\|^2 + 2\langle \bar{x}, \bar{y} \rangle + \|\bar{y}\|^2
    \]
    Using the Cauchy-Schwarz inequality, we get:
    \[
        \|\bar{x} + \bar{y}\|^2 \leq (\|\bar{x}\| + \|\bar{y}\|)^2
    \]
    Hence:
    \[
        \|\bar{x} + \bar{y}\| \leq \|\bar{x}\| + \|\bar{y}\|
    \]
    For the second inequality, we have:
    \[
        |\|\bar{x}\| - \|\bar{y}\|| = |\|\bar{x}\| - \|\bar{y}\|| = |\|\bar{x}\| - \|\bar{y}\|| \leq \|\|\bar{x}\| - \|\bar{y}\|| + \|\bar{y}\| = \|\bar{x}\| \leq \|\bar{x} - \bar{y}\|
    \]
\end{proof}

\subsection{Distance in $\mathbb{R}^n$}
\begin{definition}[Distance]
    The distance between two points $\bar{x}, \bar{y} \in \Rn$ is defined as:
    \[
        d(\bar{x}, \bar{y}) = \|\bar{x} - \bar{y}\|
    \]
\end{definition}
The distance function $d$ satisfies the following properties for all $\bar{x}, \bar{y}, \bar{z} \in \Rn$:
\begin{citemize}
    \item $d(\bar{x}, \bar{y}) = 0$ if and only if $\bar{x} = \bar{y}$
    \item $d(\bar{x}, \bar{y}) = d(\bar{y}, \bar{x})$ (symmetry)
    \item $d(\bar{x}, \bar{z}) \leq d(\bar{x}, \bar{y}) + d(\bar{y}, \bar{z})$ (triangle inequality)
\end{citemize}

\section{Open and Close Subsets of $\mathbb{R}^n$}
\begin{definition}[Open Ball]
    For all $\bar{x} \in \Rn$ and $\delta > 0$, the open ball centered at $\bar{x}$ with radius $\delta$ is defined as:
    \[
        B(\bar{x}, \delta) = \{\bar{y} \in \Rn : d(\bar{x}, \bar{y}) < \delta\} = \{\bar{y} \in \Rn : \|\bar{x} - \bar{y}\| < \delta\}
    \]
\end{definition}

\begin{definition}[Open Subset of $\mathbb{R}^n$]
    A subset $A \subseteq \Rn$ is called open if:
    \begin{citemize}
        \item $A = \emptyset$, or
        \item For all $\bar{x} \in A$, there exists $\delta > 0$ such that $B(\bar{x}, \delta) \subset A$
    \end{citemize}
\end{definition}

\begin{eg}
    Let there be an open ball in $\R$. We have:
    \[
        B(\bar{x}, \delta) = \{\bar{y} \in \R : \|\bar{x} - \bar{y}\| < \delta\} = (\bar{x} - \delta, \bar{x} + \delta) 
    \]
    Which is an open interval in $\R$. Hence, the open balls in $\R$ are the open intervals.
\end{eg}

\begin{definition}[Interior Point]
    Let $A \subseteq \Rn$ and $\bar{x} \in A$. We say that $\bar{x}$ is an interior point of $A$ if there exists $\delta > 0$ such that $B(\bar{x}, \delta) \subset A$.
\end{definition}
The set of all interior points of $A$ is called the interior of $A$ and is denoted by $\mathring{A}$. We clearly have that $\mathring{A} \subset A$ and a set A is open if and only if $A = \mathring{A}$.

\begin{eg}
    Every open ball $B(\bar{x}, \delta) \subset \Rn$ is an open set. Let $\bar{y} \in B(\bar{x}, \delta)$. We have:
    \[
        \|\bar{x} - \bar{y}\| < \delta
    \]
    Let $\epsilon = \frac{1}{2} \left(\delta - \|\bar{x} - \bar{y}\|\right) > 0$. We have:
    \[
        B(\bar{y}, \epsilon) \subset B(\bar{x}, \delta)
    \]
    Which implies that $\bar{y}$ is an interior point of $B(\bar{x}, \delta)$. Since $\bar{y}$ was arbitrary, we conclude that $B(\bar{x}, \delta)$ is an open set.
\end{eg}

\begin{eg}
    The set $E = \{\bar{x} \in \Rn : x_i > 0, \forall i = 1, \ldots, n\}$ is open. Let $\bar{y} \in E$. We have:
    \[
        B(\bar{y}, \frac{1}{2} \min(y_i)) \subset E
    \]
    Which implies that $\bar{y}$ is an interior point of $E$. Since $\bar{y}$ was arbitrary, we conclude that $E$ is an open set.
\end{eg}
The empty set $\emptyset$ and the whole space $\Rn$ are open sets.

\begin{eg}
    Let $n \geq 2$ and $E = \{\bar{x} \in \Rn : x_1 = 0, x_i > 0, i = 2, \ldots, n\}$. We have $\bar{y} = (0, y_2, y_3, \ldots, y_n) \in E$. For all $\delta > 0$, we have:
    \[
        \left(\frac{\delta}{2}, y_2, y_3, \ldots, y_n\right) \in B(\bar{y}, \delta) \quad \text{and} \quad \left(\frac{\delta}{2}, y_2, y_3, \ldots, y_n\right) \notin E
    \]
    Which implies that $\bar{y}$ is not an interior point of $E$. Since $\bar{y}$ was arbitrary, we conclude that $E$ has no interior points and is not an open set.
\end{eg}
The following properties holds for open sets in $\Rn$:
\begin{citemize}
    \item The union of any collection of open sets in $\Rn$ is an open set in $\Rn$.
    \[
        \bar{x} \in \bigcup_{i \in I} E_i \Implies \exists j : \bar{x} \in E_j \text{ open } \Implies \exists \delta > 0 : B(\bar{x}, \delta) \subset E_j \subset \bigcup_{i \in I} E_i = E
    \]
    \item The intersection of a finite number of open sets in $\Rn$ is an open set in $\Rn$.
    \[
        \bar{x} \in \bigcap_{i = 1}^n E_i \Implies \forall j : \bar{x} \in E_j \text{ open } \Implies \exists \delta_j > 0 : B(\bar{x}, \delta_j) \subset E_j
    \]
    \[
        \Implies B(\bar{x}, \min(\delta_j)) \subset E_j \forall j \Implies B(\bar{x}, \min(\delta_j)) \subset \bigcap_{i = 1}^n E_i = E
    \]
\end{citemize}
Remark that the intersection of an infinite number of open sets in $\Rn$ is not necessarily an open set in $\Rn$.
\begin{eg}
    Let $E$ be the intersection of the open sets $E_n = B(\bar{0}, \frac{1}{n})$ for $n \geq 1$. We have:
    \[
        E = \bigcap_{n = 1}^{\infty} E_n = \{\bar{0}\} \subset \Rn
    \]
    Which is not an open set in $\Rn$.
\end{eg}

\begin{definition}[Closed Subset of $\mathbb{R}^n$]
    A subset $A \subseteq \Rn$ is called closed if its complement $\complement A$ (or $A^c$) $\complement A = \Rn \setminus A$ is an open set in $\Rn$.
\end{definition}

\begin{eg}
    The complementary of an open ball $\complement B(\bar{x}, \delta) = \{\bar{y} \in \Rn : \|\bar{x} - \bar{y} \| \geq \delta\}$ is a closed set in $\Rn$ since the complement of $\complement B(\bar{x}, \delta)$ is the open ball $B(\bar{x}, \delta)$.
\end{eg}

\begin{eg}
    The point $E = \{\bar{x}\}$ is closed since the complement of $E$ is the open set $\complement E = \{\bar{y} \in \Rn : \bar{y} \neq \bar{x}\} = \{\bar{y} \in \Rn : \|\bar{y} - \bar{x} \| > 0\}$. Thus we have:
    \[
        \forall \bar{y} \in \complement E, \quad \exists \delta = \frac{1}{2} \|\bar{y} - \bar{x}\| > 0 \quad : \quad B(\bar{y}, \delta) \subset \complement E
    \]
\end{eg}

\begin{eg}
    The hyperplane $E = \{\bar{x} \in \Rn : x_1 = 0\}$ is closed since the complement of $E$ is the open set $\complement E = \{\bar{y} \in \Rn : y_1 \neq 0\} = \{\bar{y} \in \Rn : |y_1| > 0\}$. Thus we have:
    \[
        \forall \bar{y} \in \complement E, \quad \exists \delta = \frac{1}{2} |y_1| > 0 \quad : \quad B(\bar{y}, \delta) \subset \complement E
    \]
\end{eg}

\begin{eg}
    The set $E = \{\bar{x} \in \Rn : x_1 = 0, x_i > 0 \ \forall i \geq 2\}$ ($n \geq 2$) can be defined as:
    \[
        E = \{\bar{x} : x_1 = 0\} \cap \{\bar{x} : x_2 > 0\} \cap \ldots \cap \{\bar{x} : x_n > 0\}
    \]
    Hence the complement of $E$ can be defined as:
    \[
        \complement E = \{\bar{x} : x_1 \neq 0\} \cup \{\bar{x} : x_2 \leq 0\} \cup \ldots \cup \{\bar{x} : x_n \leq 0\}
    \]
    We can take $\bar{y} = (0,0,y_3 > 0, \ldots, y_n > 0) \in \complement E$ and thus for all $\delta > 0$, we have:
    \[
        (0, \frac{\delta}{2}, y_3, \ldots, y_n) \in B(\bar{y}, \delta) \quad \text{and} \quad (0, \frac{\delta}{2}, y_3, \ldots, y_n) \notin \complement E
    \]
    Which implies that $\bar{y}$ is not an interior point of $\complement E$. Since $\bar{y}$ was arbitrary, we conclude that $\complement E$ is not an open set in $\Rn$ and thus $E$ is not a closed set in $\Rn$ ($E$ is not open and not close).
\end{eg}
Remark that the only open and closed sets in $\Rn$ are the empty set $\emptyset$ and the whole space $\Rn$ since the complement of $\emptyset$ is $\Rn$ and the complement of $\Rn$ is $\emptyset$. We also have the following properties:
\begin{citemize}
    \item The intersection of any collection of closed sets in $\Rn$ is a closed set in $\Rn$ ($\bigcap_{i \in I} E_i$).
    \[
        \complement E = \bigcup_{i \in I} \complement E_i \text{ open } \Implies E = \bigcap_{i \in I} E_i \text{ closed }
    \]
    \item The union of a finite number of closed sets in $\Rn$ is a closed set in $\Rn$ ($\bigcup_{i = 1}^n E_i$).
    \[
        \complement E = \bigcap_{i = 1}^n \complement E_i \text{ open } \Implies E = \bigcup_{i = 1}^n E_i \text{ closed }
    \]
\end{citemize}

\section{The closure and the boundary of a subset of $\mathbb{R}^n$}
\begin{definition}[Closure]
    Let $E$ be a non-empty subset of $\Rn$, the closure of $E$ is defined as the intersection of all closed sets in $\Rn$ that contain $E$, denoted by $\overline{E}$.
\end{definition}
Remark that $E \subset \Rn$ is closed if and only if $E = \overline{E}$.

\begin{definition}[Boundary]
    Let $E$ be a non-empty and proper subset of $\Rn$. Given a point $\bar{x} \in \Rn$, we say that $\bar{x}$ is a boundary point of $E$ if for all open ball $B(\bar{x}, \delta)$ centered at $\bar{x}$ contains at least one point of $E$ and at least one point of $\complement E$.
\end{definition}
Remark that the set of all boundary points of $E$ is called the boundary of $E$ and is denoted by $\partial E$.

\begin{eg}
    Let $E = \{\bar{x} \in \Rn : x_i > 0, \ i = 1, \ldots, n\}$ an open set in $\Rn$. We have:
    \[
        \overline{E} = \{\bar{x} \in \Rn : x_i \geq 0, \ i = 1, \ldots, n\}
    \]
    and
    \[
        \partial E = \{\bar{x} \in \Rn : \exists i : x_i = 0 \text{ and } \forall j \neq i \implies x_j \geq 0\}
    \]
\end{eg}
Let $E \subset \Rn$ be a non-empty subset of $\Rn$. We have the following properties:
\begin{citemize}
    \item $\partial E \cap \mathring{E} = \emptyset$ (where $\mathring{E}$ is the interior of $E$)
    \item $\mathring{E} \cup \partial E = \bar{E}$ (since $\mathring{E} \cup \partial E$ is closed and contains $E$, it's the smallest closed set that contains $E$)
    \item $\partial E = \bar{E} - \mathring{E} = \bar{E} \cap \complement \mathring{E}$ and thus $\partial E$ is closed.
    \item $\partial \emptyset = \emptyset$ and $\partial \Rn = \emptyset$
\end{citemize}

\section{Sequence of Elements of $\mathbb{R}^n$ and Topology of $\mathbb{R}^n$}
\begin{definition}[Sequence of Elements of $\mathbb{R}^n$]
    A sequence of elements of $\Rn$ is a function $f : \N \to \Rn$ such that:
    \[  
        f: k \to \bar{x}_k = (x_{k,1}, x_{k,2}, \ldots, x_{k,n} ) \in \Rn
    \]
    We denote the sequence by $\{\bar{x}_k\}_{k = 0}^{\infty}$.
\end{definition}

\begin{definition}[Convergence of a Sequence of Elements of $\mathbb{R}^n$]
    Let $\{\bar{x}_k\}_{k = 0}^{\infty}$ be a sequence of elements of $\Rn$ and $\bar{x} \in \Rn$. We say that $\{\bar{x}_k\}_{k = 0}^{\infty}$ converges to $\bar{x}$ if for all $\epsilon > 0$, there exists $N \in \N$ such that for all $k \geq N$, we have:
    \[
        d(\bar{x}_k, \bar{x}) = \| \bar{x}_k - \bar{x} \| \leq \epsilon
    \]
    We denote this by $\lim_{k \to \infty} \bar{x}_k = \bar{x}$ or $\bar{x}_k \to \bar{x}$ as $k \to \infty$.
\end{definition}
Let $\bar{x} = (x_1, \ldots, x_n)$ and $\bar{x}_k = (x_{1,k}, \ldots, x_{n,k})$ then $\lim_{k \to \infty} \bar{x}_k = \bar{x}$ if and only if for all $j \in \{1, \ldots, n\}$, $\lim_{k \to \infty} x_{j,k} = x_j$. Indeed:
\[
    \| \bar{x}_k - \bar{x} \| = \sqrt{\sum_{j = 1}^{n} (x_{j,k} - x_j)^2}
\]
We have:
\begin{citemize}
    \item If $\| \bar{x}_k - \bar{x} \| \leq \epsilon$, then for each $j$, we have $|x_{j,k} - x_j| \leq \epsilon$.
    \item If $|x_{j,k} - x_j | \leq \epsilon$ for all $j$, then we have $\| \bar{x}_k - \bar{x} \| \leq \sqrt{n}\epsilon$.
\end{citemize}

\begin{definition}[Bounded Sequence of Elements of $\mathbb{R}^n$]
    A sequence $\{\bar{x}_k\}_{k = 0}^{\infty}$ of elements of $\Rn$ is called bounded if it's contained in a closed ball $B(\bar{0}, M)$ for a given $M > 0$, i.e. if there exists $M > 0$ such that for all $k \in \N$, we have $\|\bar{x}_k\| \leq M$.
\end{definition}
We have the following properties of convergence in $\Rn$:
\begin{citemize}
    \item The limit of a sequence in $\Rn$ is unique.
    \item Every convergent sequence in $\Rn$ is bounded.
\end{citemize}

\begin{theorem}[Bolzano-Weierstrass Theorem]
    Every bounded sequence in $\Rn$ has a convergent subsequence.
\end{theorem}

\begin{theorem}
    Let $E \subset \Rn$ be a non-empty subset. $E$ is closed if and only if for every sequence $\{\bar{x}_k\} \subset E$ that converges to $\bar{x} \in \Rn$ has its limit $\bar{x}$ in $E$.
\end{theorem}
\begin{proof}
    % Let's prove both directions of the equivalence:  \\
    ($\Rightarrow$) By contradiction, let $E$ be closed and let $\bar{x} = \lim_{k \to \infty} \bar{x}_k$ with $\bar{x}_k \in E$. Suppose that $\bar{x} \notin E$. We have that $\bar{x} \in \complement E$ which is open. Hence there exits $\delta > 0$ such that $B(\bar{x}, \delta) \subset \complement E$. This implies that $\{\bar{x}_k\} \cap B(\bar{x}, \delta) = \emptyset$, which contradicts the fact that $\bar{x}_k$ converges to $\bar{x}$. For example, we can take $k_0 \in \N$ such that for all $k \geq k_0$ we have $\bar{x}_k \in \overline{B(\bar{x}, \delta / 2)} \subset B(\bar{x}, \delta)$. \\
    ($\Leftarrow$) By contraposition, let $E$ be not closed (then $\complement E$ is not open). There exists $\bar{y} \in \complement E$ such that for all $k \in \N^*$, $B(\bar{y}, 1/k) \cap E \neq \emptyset$. Let $y_k \in B(\bar{y}, 1/k) \cap E$, we can construct a sequence $\{\bar{y}_k\} \subset E$ such that $\bar{y}_k \in B(\bar{y}, 1/k)$. Then $\lim_{k \to \infty} \bar{y}_k = \bar{y}$. Since $\bar{y} \notin E$, we conclude that there exists a sequence $\{\bar{y}_k\} \subset E$ that converges to $\bar{y} \notin E$.
\end{proof}
Remark that we can build the closure of a set $E$ by taking the limits of all convergent sequences in $E$.

\begin{definition}[Compact Subset of $\mathbb{R}^n$]
    A non-empty subset $E \subset \Rn$ is called compact if it is closed and bounded.
\end{definition}

\begin{eg}
    A closed ball $\overline{B(\bar{x}, \delta)}$ is compact since it's closed and bounded. Since we have:
    \[
        \overline{B(\bar{x}, \delta)} \subset \overline{B(\bar{0}, \| \bar{x} \| + \delta)} 
    \]
\end{eg}

\begin{eg}
    The set $E = \{\bar{x} \in \Rn : n \geq 2, \ x_1 = 0\}$ is closed but not bounded and thus not compact. We can take:
    \[
        \{\bar{a}_k = (0, k, 0, \ldots)\}_{k \in \N} \subset E
    \]
    but the normes of the elements of this sequence are not bounded since $\|\bar{a}_k\| = k$ for all $k \in \N$.
\end{eg}

\begin{eg}
    An open ball $B(\bar{x}, \delta)$ is not compact since it's not closed, but it is bounded since:
    \[
        B(\bar{x}, \delta) \subset \overline{B(\bar{0}, \| \bar{x} \| + \delta)}
    \]
\end{eg}

\subsection{Heine-Borel Theorem}
\begin{theorem}[Heine-Borel Theorem]
    A non-empty subset $E \subset \Rn$ is compact if and only if every collection of open sets in $\Rn$ that covers $E$:
    \[
        E \subset \bigcup_{i \in I} E_i \text{ with } E_i \text{ open in } \Rn
    \]
    where $I$ can be an infinite (countable or uncountable) index set, has a finite sub-collection that also covers $E$:
    \[
        E \subset \bigcup_{i = 1}^n E_{j_i} \text{ with } j_i \in I \text{ for all } i = 1, \ldots, n
    \]
\end{theorem}

\begin{eg}
    A line $E \subset \R^2$ is not compact since it's not bounded. We can take the following collection of open sets that covers $E$:
    \[
        E \subset \bigcup_{n \in \Z} B\left(n, \frac{2}{3}\right)
    \]
    but it is impossible to find a finite sub-collection of this collection that covers $E$ thus $E$ is not compact.
\end{eg}

\begin{eg}
    The open interval $E = (0,1) \subset \R$ is not compact since it's not closed. We can take the following collection of open sets that covers $E$:
    \[
        E \subset \bigcup_{n \in \N} \left(0, \frac{n}{n + 1}\right)
    \]
    but here again, it is impossible to find a finite sub-collection of this collection that covers $E$ thus $E$ is not compact.
\end{eg}

\begin{eg}
    Let $S = \{(x,y) \in \R^2 : \sqrt{y \cdot \ln x} < 1\}$. Is $S$ closed or opened? Bounded or unbounded? We have:
    \begin{citemize}
        \item We have $\ln x$ in the formula, thus $x > 0$.
        \item Let $y = 0$, we have $\{y = 0, x > 0\} \subset S$ and $\{x = 1, y \in \R\} \subset S$.
        \item Let $y > 0$, we have:
        \[
            y \ln x \geq 0 \Implies \ln x \geq 0 \Implies x \geq 1
        \]
        Thus:
        \[
            y \ln x < 1 \Implies y < \frac{1}{\ln x}
        \]
        Hence:
        \[
            \left\{y > 0, x > 1, y < \frac{1}{\ln x}\right\} \subset S
        \]
        \item Let $y < 0$, we have:
        \[
            y \ln x \geq 0 \Implies \ln x \leq 0 \Implies 0 < x \leq 1
        \]
        Thus:
        \[
            y \ln x < 1 \Implies y > \frac{1}{\ln x}
        \]
        Hence:
        \[
            \left\{y < 0, 0 < x < 1, y > \frac{1}{\ln x}\right\} \subset S
        \]
    \end{citemize}
    Graphically, we have:
    % TODO: make graph here
    \begin{center}
        \begin{tikzpicture}
            % \begin{axis}[
            %     axis lines = middle,
            %     xlabel = $x$,
            %     ylabel = $y$,
            %     xtick={0,1},
            %     ytick={0},
            %     ymin=-2, ymax=2,
            %     xmin=-1, xmax=3,
            %     legend pos=outer north east
            % ]
            % \addplot [
            %     domain=0.01:3, 
            %     samples=100, 
            %     color=blue,
            % ]
            % {1/ln(x)};
            % \addlegendentry{$y = \frac{1}{\ln x}$}
            
            % \addplot [
            %     domain=0.01:3, 
            %     samples=100, 
            %     color=red,
            % ]
            % {-1/ln(x)};
            % \addlegendentry{$y = -\frac{1}{\ln x}$}
            
            % \end{axis}
        \end{tikzpicture}
    \end{center}
    We can see that:
    \begin{citemize}
        \item $S$ is not open since any point on the line $x = 1$ will not have an open ball centered at it that is contained in $S$.
        \item $\complement S$ is not open since any point on the curve $y = \frac{1}{\ln x}$ will not have an open ball centered at it that is contained in $\complement S$ thus $S$ is not closed.
        \item $S$ is clearly unbounded since, for example, for $x = 1$ we have $y \in \R$ and thus $S$ contains the line $\{x = 1, y \in \R\}$ which is unbounded.
    \end{citemize}
\end{eg}